\scanpage{122}
$X$ Schéma \unsure{compact}, [i.e. \unsure{quasi-compact} et \inserted{\uline{\unsure{topologiquement}}} \uuline{séparé}]

(i) $X$ est tot. discontinu, \deleted{i.e. si $x,y\in X$, $x\ne y$, alors il existe une \unsure{partie ouverte et fermée} de $X$ qui contient $x$ et non $y$. En effet, \unsure{soient} $U,V$ des ouverts, $U\ni x$, $V\ni y$, $U\cap V=\varnothing$.}
i.e. tt $x\in X$ a un syst. fond. de voisinages qui sont à la fois \uline{ouverts} et \uline{fermés}.

[En effet, il y a un syst. fond. de voisinages ouverts quasi-compacts]

(ii) $X$ \unsure{est} affine. En effet, \deleted{\unsure{prenons} \illegible{6} \illegible{3} \unsure{affines} \illegible{5}} \inserted{on recouvre $X$ par des ouverts affines en \unsure{n\textsuperscript{bre}} fini, soit $X_i$ $(1\leq i\leq n)$}. \deleted{\unsure{les}.} Ils sont ouverts \unsure{et} fermés. Soit
\[
X'_i=X_i-X_i\cap\bigcup_{j<i}X_j,
\]
Alors $X'_i$ est une partie ouverte et fermée de $X_i$, donc affine puisque $X_i$ l'est (\unsure{en vertu} p.ex. du \unsure{critère élém. pour} \unsure{immersions} \unsure{fermés|fermées}). Donc $X$ est \deleted{\unsure{et}} somme des $X'_i$, il est affine.

(iii) Les parties constructibles de $X$ sont des parties ouvertes et fermées. [\deleted{Donc} \inserted{\unsure{En eff.}} car $\mathrm{Cons}\,X\to X$ est un \deleted{\unsure{homéomorphisme}} \unsure{hom|bon}] Donc
\margin{121}
